\chapter{Практика}
\label{ch:chap2}

% \begin{figure}[H]
%     \centering
%     \includegraphics{carrier.png}
%     \caption{Несущие колебания}
% \end{figure}

\section{Построение случайного процесса}
Построим плотность распределения для случайного процесса с нормальным распределением. Проследим, как меняется плотность распределения
при изменении мат.ожидания и дисперсии.

\begin{figure}[h]
\centering

\begin{subfigure}{0.75\textwidth}
    \includegraphics[width=\linewidth]{001.png}
    \caption{mu=0, sigma=0.1}
\end{subfigure}
\hfill
\begin{subfigure}{0.75\textwidth}
    \includegraphics[width=\linewidth]{515.png}
    \caption{mu=5, sigma=1.5}
\end{subfigure}

\caption{Примеры плотности распределения нормального распределения}
\end{figure}

Можно заметить, что при мат.ожидание mu влияет на центр плотности распределения (смещение по оси Х), а дисперсия sigma отвечает
за ширину плотности распределения или за разброс.

\section{Выборка из нормального распределения и гистограмма плотности распределения}
Сделаем выборку из нормального распределения и построим гистограмму плотности распределения с помощью plt.hist()

\begin{figure}[H]
    \centering
    \includegraphics{2.png}
    \caption{mu=7, sigma=25}
\end{figure}

Таким образом, если у нас есть случайная величина, то мы можем построить гистограмму и проанализировать эту величину.

\section{Числовые характеристики случайной величины}
Вычислим дисперсию и мат.ожидание по полученной СВ:

$$\boxed{mu = \frac{1}{N}\sum_{n=1}^{N}x[n]}$$
$$\boxed{D = \frac{1}{N}\sum_{n=1}^{N}x^2[n] - mu^2}$$

Результат:

\begin{figure}[H]
    \centering
    \includegraphics{3.png}
    \caption{Эмирические числовые характеристики СВ}
\end{figure}

Я задавал СВ с нормальным распределением с mu=7, sigma=25. Таким образом мы можем вычислять числовые характеристики СВ, имея только
саму реализацию СВ.



\endinput 